\chapter{Prototype 1: Infrastructure and CI/CD Scaffolding}
\label{app:proto1}

This appendix outlines the architecture and configuration of the first
prototype developed during this project (\texttt{java-ee-app}). Unlike the
second prototype (detailed in Appendix~\ref{app:syncspace}), which focused on
Jakarta EE feature exploration, this initial application was built with a
strictly constrained scope: to serve as a proof-of-concept for the underlying
infrastructure. Its primary objective was to validate the integration
between the Payara 7 application server, Arquillian integration testing,
and the GitLab \ac{CI}/\ac{CD} pipeline before introducing complex domain
logic.

\section{Technology Stack}
\label{app:proto1-stack}

The prototype establishes the baseline toolchain that was later adapted for both
the SyncSpace prototype and the production \ac{EPOS} \ac{API} refactoring.

\begin{table}[H]
	\centering
	\renewcommand{\arraystretch}{1.2}
	\begin{tabular}{ll}
		\toprule \textbf{Concern} & \textbf{Technology}                                   \\
		\midrule Platform         & Jakarta EE 11 (\texttt{jakarta.jakartaee-api 11.0.0}) \\
		Language                  & Java 17                                               \\
		App Server                & Payara 7 (\texttt{7.2026.x})                          \\
		Build                     & Maven 3.9.6 (\ac{WAR} packaging)                      \\
		Database                  & PostgreSQL 16 (Alpine Docker image)                   \\
		\ac{REST} Layer           & Jakarta \ac{REST}ful Web Services (\ac{JAX-RS})       \\
		Testing                   & Arquillian 1.10 + JUnit 5 (Payara Managed container)  \\
		\ac{CI}/\ac{CD}           & GitLab CI                                             \\
		\bottomrule
	\end{tabular}
	\caption{Prototype 1 baseline technology stack.}
\end{table}

\section{Architecture and Endpoints}
\label{app:proto1-arch}

To verify basic deployment functionality and content negotiation, the
application exposes a headless \ac{REST} \ac{API} utilizing two independent
\ac{JAX-RS} \texttt{Application} subclasses. This design deliberately separates
a legacy plain-text route from a versioned \ac{JSON} route:

\begin{itemize}
	\item \textbf{\texttt{HelloApplication} (\texttt{/api}):} Registers a
		simple \texttt{HelloResource} returning a \texttt{text/plain} fixed
		string (\texttt{"Hello, World!"}).

	\item \textbf{\texttt{JsonApplication} (\texttt{/api/v2}):} Registers a
		\texttt{JsonResource} returning \texttt{application/json}. It serializes
		a minimal \texttt{User} \ac{POJO} containing an auto-generated UUID, a
		name, and an age.
\end{itemize}

Because the scope of this prototype was restricted to infrastructure validation,
cross-cutting concerns such as security, input validation, and \ac{CDI}-managed
services were intentionally left unimplemented.

\section{Persistence Readiness}
\label{app:proto1-persistence}

While the application does not actively execute database queries, the scaffolding
for persistence is fully established to validate the environment setup. A
\texttt{persistence.xml} file exists defining a default persistence unit,
and a PostgreSQL 16 instance (\texttt{testdb}) is provisioned locally via Docker
Compose alongside the Payara server. The domain model (the \texttt{User}
class) remains a plain Java object rather than a \ac{JPA} \texttt{@Entity},
leaving the final persistence wiring as a future roadmap item for the
architecture.

\section{Testing Framework}
\label{app:proto1-testing}

The most critical deliverable of this prototype is the integration testing harness.
The project utilizes Arquillian configured for a Payara Managed container, combined
with JUnit 5.

The testing lifecycle is anchored by an \texttt{AbstractBaseIT} class. This
base class uses ShrinkWrap to programmatically construct a micro-\ac{WAR} (\texttt{test-api.war})
from the application's source packages. It handles the setup and teardown of
a \ac{JAX-RS} \texttt{Client} per test and injects the dynamic deployment base
URL via the \texttt{@ArquillianResource} annotation. Specific integration tests,
such as \texttt{JsonApplicationIT}, inherit from this base class to issue
actual \ac{HTTP} requests against the running container and assert successful
\texttt{200 OK} \ac{JSON} responses.

\section{CI/CD Pipeline}
\label{app:proto1-pipeline}

The application's lifecycle is automated via a \texttt{.gitlab-ci.yml} pipeline
executing on a self-managed runner. The pipeline successfully validated the
two-stage execution model later adopted by the broader project:

\begin{enumerate}
	\item \textbf{Build Stage:} Executes \texttt{mvn clean package -DskipTests}
		within a Temurin JDK Docker image, caching dependencies in the global
		\texttt{.m2} repository, and publishes the resulting \ac{WAR} as a pipeline
		artifact.

	\item \textbf{Test Stage:} Boots the Payara service container, waits for readiness,
		and executes the Failsafe integration tests via \texttt{mvn verify}.
		The results are captured as JUnit \ac{XML} artifacts and ingested into GitLab's
		native test tracking interface.
\end{enumerate}

By successfully executing this pipeline, the prototype confirmed that the
laboratory's self-managed runners could reliably compile Jakarta EE
applications, orchestrate containerized application servers, and execute
robust integration tests over \ac{HTTP}.